% ============================================================================
% CAPÍTULO 6 — EXTENSIONES Y TRABAJO FUTURO
% Objetivo: dejar planteadas, CON PROPIEDAD, las extensiones tratables.
% Contenido (según lo discutido, ordenado por costo real):
%   · Carga (Reissner-Nordström): cambio de f_+(R); la forma cerrada sobrevive
%     con coeficiente (b^2-Q^2)/b. Costo bajo. Refs: Kuchar 1968; Eid 2015.
%   · Rotación lenta O(Omega): NO rompe R(tau) a primer orden; agrega una
%     ecuación para omega(r). Arrastre perfecto en r_s (resultado on-theme
%     black shell). SECCIÓN REAL (decisión de Santiago). Refs: Hartle 1967;
%     Brill & Cohen 1966.
%   · Presión (Núñez): m(R) deja de ser constante; permite bounce y black
%     shell estático de verdad. Refs: Núñez 1996, 1997.
% NOTA: Kerr completo NO es una extensión — rompe el método (no hay Birkhoff).
%   Mencionar por qué queda fuera, para cerrar esa puerta con argumento.
% ============================================================================
\chapter{Extensions and future work}
\label{cap:extensiones}

\section{Cáscara cargada (Reissner--Nordström)}
% TODO: f_+ = 1 - 2M/R + Q^2/R^2 ; forma cerrada con (b^2-Q^2)/b (VERIFICAR
%   contra Kuchar 1968 / Eid 2015 el tratamiento de la autocarga).

\section{Rotación lenta y arrastre de marcos}
\label{sec:rotacion}
% TODO: expansión de Hartle O(Omega); ruptura confinada a g_{t phi};
%   omega_+(r)=2J/r^3 (Lense-Thirring); interior arrastrado (Brill-Cohen);
%   arrastre perfecto omega_- -> Omega cuando R -> r_s.

\section{Presión y black shell estático}
% TODO: generalización de m(R); conexión con bounce (Núñez).

\section{Por qué Kerr queda fuera}
% TODO: sin simetría esférica no hay Birkhoff; el shell deja de ser R(tau).
